% !TeX root = ../nc-cubic-non-as.tex
\subsection{巡回関係の三次syzygyと行列式rank二の認識}
\label{subsec:ou-cyclic-determinants}

前節の三ブロック関係代数$\Lambda$を考える。
係数は$c_{ijk}$、矢印は$A_{ij}:U_i\to V_j$、
$B_{jk}:V_j\to W_k$とする。
添字$i,j,k$はそれぞれ$0,1,2$を動く。
この節では$\Perf\Lambda$上の導来helixがvery strongであると仮定する。
特に次の周期の第一ブロックは$sU_i$である
（$s=S_\Lambda^{-1}[2]$）。
まず曲線表示を用いずに巡回した三ブロック表示を記述する。

\begin{lemma}[巡回表示と三次syzygy]
\label{lem:ou-cyclic-cubic-syzygy}
逆Serreの周期をまとめたhelix代数には、$W_k$から$sU_i$への
矢印$C_{ki}$があり、三種類の二次関係は
\[
 \sum_j c_{ijk}B_{jk}A_{ij}=0,\qquad
 \sum_k c_{ijk}C_{ki}B_{jk}=0,\qquad
 \sum_i c_{ijk}A_{ij}C_{ki}=0
\]
で与えられる（最後の式では次の周期の$A_{ij}$を意味する）。
従って三ブロックの巡回は係数テンソルの巡回置換で実現される。

各方向の二次関係の係数ベクトルは非零である。
さらに$i,i'$に対して、六未知数$\alpha_k,\beta_j$の連立方程式
\begin{equation}\label{eq:ou-cyclic-syzygy}
 \alpha_k c_{ijk}=\beta_j c_{i'jk}
 \qquad(0\leq j,k\leq2)
\end{equation}
の解空間は$i=i'$なら一次元、$i\ne i'$なら零である。
添字を巡回置換した同じ条件が成り立つ。
\end{lemma}
\begin{proof}
主稿の定理\ref{thm:ou-weak-cubic}の箙型の確認と同じ逆Serre
bimoduleの計算を、一般の三ブロック代数に適用する。
$\Theta=R\operatorname{Hom}_{\Lambda^e}(\Lambda,\Lambda^e)$と置くと
$s$は$-\otimes_\Lambda^{\mathbf L}\Theta[2]$である。
$\Lambda$の頂点・矢・二次関係に対応する長さ二の最小bimodule分解は、
normalized bar resolutionの同型部分の消去で得られる。
very strong性により$\Theta[2]$の全正tensor冪は次数零に集中するので、
周期をまとめた代数は$T_\Lambda H^0(\Theta[2])$である。
双対分解の最終余核の生成元は九つの関係を逆に向けた$C_{ki}$であり、
微分を双対化すると、元の二次関係の各$A_{ij},B_{jk}$を
一つ取り除いて逆矢印を掛けた二種類の関係を得る。
例えば$B_{jk}A_{ij}$の微分は
$[B_{jk}]A_{ij}+B_{jk}[A_{ij}]$である。
これが上の三種類の巡回関係を与える。
従ってこの表示は
AOU\cite[(4.7), (4.10), Proposition~5.1]{AOU2024}の
potentialの表示と一致する。ここでは任意のpotentialが
AS正則になるとは仮定しない。

元の射影基底のEuler行列を$D$とすると、
補題\ref{lem:ou-residual-helix-growth}の計算により
\[
 S_0=D^{-1}D^{\mathsf t},\qquad
 S_0^{\mathsf t}D=
 \begin{pmatrix}
 I_3+3J_3&5J_3&7J_3\\
 2J_3&I_3+3J_3&5J_3\\
 J_3&2J_3&I_3+3J_3
 \end{pmatrix}.
\]
特にvery strong性から、巡回した各二段経路のHom次元は二で、
三段経路のHom次元は$3+\delta_{ii'}$である。
二段経路の三本のうち関係が一つ必要なので、
全方向で係数ベクトルは非零である。
$U_i$から$sU_{i'}$への三段経路は九本。
最初の二段と最後の二段に由来する関係部分空間は
それぞれ三次元であり、その交わりの次元は
\[
 (3+3)-(9-(3+\delta_{ii'}))=\delta_{ii'}.
\]
二つの部分空間の元が一致する条件は、各経由組$(j,k)$について
\eqref{eq:ou-cyclic-syzygy}そのものである。
巡回したfoundationについても同じEuler行列と議論が適用できる。
\end{proof}

\begin{corollary}[係数の台に対する必要条件]
\label{cor:ou-cyclic-support-obstruction}
各方向の固定添字の$3\times3$係数行列について、その非零成分が
定める三行・三列の二部graphは連結である。
異なる二つの固定添字の行列$C,C'$については、
両方が非零の成分を辺とし、一方だけが非零の成分により
対応する変数を零とする二部graphを考える。
零変数を含まない連結成分があるなら、その成分内の
比$c_{ijk}/c_{i'jk}$は辺上の連立方程式と整合してはならない。
特に零変数を含まないtree成分（孤立頂点を含む）があれば
very strong性に反する。

係数が全て非零の場合には、この三次syzygyの条件は、任意の方向の
相異なる二sliceについてその成分ごとの比の$3\times3$行列の
rankが一でないことと同値である。
\end{corollary}
\begin{proof}
同じsliceでは非零辺に沿って$\alpha_k=\beta_j$となるので、
解空間の次元は連結成分数である。
異なるsliceでは、一方の係数だけが非零なら対応する変数は零。
両方が非零の辺では一方の変数が他方の非零定数倍となる。
従って零は連結成分全体に伝播する。
零の強制がない成分では、辺の比を辿ることで全変数が
一つの変数の倍となる。その変数が非零にできる必要十分条件は
全てのcycleに沿う比の積が整合することである。
treeにはその障害がない。
最後の場合は完全二部graphなので、非零解の存在は
成分ごとの比が行・列スカラーの積となることと同値である。
補題\ref{lem:ou-cyclic-cubic-syzygy}を適用して結論を得る。
\end{proof}

以後、この節では$c_{ijk}\ne0$を全ての添字について仮定する。
前節の行列式座標を、偶置換に対応する$x_1,x_2,x_3$、
奇置換に対応する$y_1,y_2,y_3$と書き換え、
\[
 T=V(F)\subset\mathbb P^5,\qquad
 F=x_1x_2x_3-y_1y_2y_3
\]
とする（前節の$z$座標の並べ替えに過ぎない）。
各行列式線形式$\ell_k$の係数ベクトルを$\lambda_k$と書く。
その偶側成分は置換に沿う三係数の積、奇側成分はその負である。
従って各$\lambda_k$は全座標非零で、
\[
 T^\vee_{\rm bin}=V(G)\subset(\mathbb P^5)^*,\qquad
 G=\lambda_{x_1}\lambda_{x_2}\lambda_{x_3}
    +\lambda_{y_1}\lambda_{y_2}\lambda_{y_3}
\]
上に属する。以下のGauss写像の計算から、$T^\vee_{\rm bin}$は
実際に$T$の射影双対超曲面であるが、
双対性の一般理論は必要ない。

\begin{lemma}[全係数非零の場合の行列式rank]
\label{lem:ou-dense-determinant-rank}
三つの$\ell_k$は非零で任意の二つは比例しない。
それらのrankが二の場合、三係数点を含む直線
$L=\mathbb P\langle\lambda_0,\lambda_1,\lambda_2\rangle$は
$T^\vee_{\rm bin}$に含まれず、
$L\cap T^\vee_{\rm bin}$は相異なる三つの全座標非零点
$[\lambda_k]$だけからなる被約スキームである。
rankが三の場合、共通零平面$\Pi$は$T$に含まれない。
これらは三つの巡回表示の各々について成立する。
\end{lemma}
\begin{proof}
$\ell_k$と$\ell_{k'}$が比例するなら、二つの固定$k$ sliceの
成分ごとの比の行列について、全六置換積が同じである。
任意の二行二列の長方形を、その外の一辺を共有する
偶・奇置換で比較すると、比の行列の全二次minorは零。
系\ref{cor:ou-cyclic-support-obstruction}に反する。
従ってrank一は起こらない。

rank二で$L\subset T^\vee_{\rm bin}$と仮定する。
$L\simeq\mathbb P^1$上の各座標の制限は非零一次形式。
$G|_L=0$なので多項式環の一意分解性から、偶側の三一次形式を
奇側の三一次形式と一対一に対応させ、それぞれ比例させられる。
偶置換と奇置換は各組で一つの辺$(i,j)$を共有する。
三偶置換と三奇置換の任意のbijectionについて、三つの共有辺は
三行三列のある一行または一列の三辺をなす。
これは$\mathfrak S_3$の六つのbijectionを列挙すれば直ちに分かる
（再現コードでも検証する）。
従って例えばある一行の三辺が共有される場合には、
残る二行$i,i'$の全ての二列の交叉比
\[
 \frac{c_{ijk}c_{i'j'k}}{c_{ij'k}c_{i'jk}}
\]
が$k$に依らない（置換積の比から共有辺の係数を消去する）。
これは$i,i'$ sliceの成分ごとの比が$j,k$の積に分離すること、
すなわちrank一であることと同値である。
一列が共有される場合も添字を入れ替えた同じ障害となる。
従って$L$は$T^\vee_{\rm bin}$に含まれない。
次数三の方程式の直線への制限は非零で、既に相異なる三点を持つので、
これらは単根であり他の交点は存在しない。

rank三で$\Pi\subset T$と仮定する。
$\Pi$上の六座標が全て非零一次形式なら、
$F|_\Pi=0$から同様に偶・奇の一次形式を一対一に比例させられる。
これらは平面の一次形式を張るので、三組の代表一次形式は一次独立。
各$\ell_k$は$\Pi$上で零だから、対になった二係数の比が
全$k$で同じである。先の共有辺と交叉比の議論で矛盾する。

$\Pi$上である座標が恒等的に零なら、$F|_\Pi=0$により
偶側と奇側に少なくとも一つずつ恒等的に零の座標がある。
これらの座標を$x_1,y_1$とする。
$\ell_k$の張る三次元ベクトル空間にはその二つの座標汎関数が属する。
従ってこの空間の元は$u\lambda_0+v e_{x_1}+w e_{y_1}$と書ける
（$\lambda_0$は全座標非零）。
残り四座標は$u$の非零定数倍なので、$G$のこの空間への制限は
$u^2$と非零一次形式の積である。
三つの$\lambda_k$は全座標非零なので$u\ne0$、
従ってすべてこの一次形式の零空間に属する。
これは三式の一次独立性に反する。
巡回性は補題\ref{lem:ou-cyclic-cubic-syzygy}による。
\end{proof}

\begin{theorem}[rank二の行列式切断は滑らかな三次曲面]
\label{thm:ou-dense-rank-two-smooth-surface}
全係数非零、導来helixのvery strong性の下で、
$\ell_0,\ell_1,\ell_2$のrankを二とする。
このとき
\[
 X=T\cap L^\perp\subset L^\perp\simeq\mathbb P^3
\]
は滑らかな三次曲面である。ここで$L$は前補題の双対直線、
$L^\perp$は共通零空間である。
従って前節の滑らかな平面切断からの復元が適用でき、
標準noetherian OU helix比較と内在的条件の下で
標準AS四次元三次商表示を得る。
\end{theorem}
\begin{proof}
補題\ref{lem:ou-dense-determinant-rank}により$L$と
$T^\vee_{\rm bin}$の交点は三つの全座標非零の単純交点である。
特に$L$には、$G=0$かつ座標が零となる点は存在しない。
$L^\perp$は四次元ベクトル空間を射影化したものであり、
$F$の制限は恒等的に零ではない。
実際、六座標が全て非零一次形式として制限されるなら、
恒等的な零は一意分解性により三組の比例する一次形式を与え、
四次元の双対空間を張ることに反する。
ある座標の制限が零なら、反対側の座標の一つも零となり、
$L$が二つの座標汎関数を含むことになるが、
これは$L\cap T^\vee_{\rm bin}$に座標点が現れるため不可能。
よって$X$は三次超曲面である。

$z\in X$が全座標非零なら、
$P=x_1x_2x_3=y_1y_2y_3\ne0$とおくと
\[
 \lambda=\nabla F(z)=
 (P/x_1,P/x_2,P/x_3,-P/y_1,-P/y_2,-P/y_3).
\]
この$\lambda$は$G=0$を満たし、
$\nabla G(\lambda)=Pz$である。
$z$が$X$の特異点なら$\nabla F(z)$は$L$に属する。
しかも$z\in L^\perp$なので$\nabla G(\lambda)$は$L$を消し、
$L$は$\lambda$で$T^\vee_{\rm bin}$に接することになる。
これは三交点が単純であることに反する。

$z$が座標零の境界にあれば、偶側・奇側にそれぞれ零がある。
$\nabla F(z)\ne0$なら、このベクトルの台は
偶側・奇側に高々一つずつなので、$G=0$を満たす境界点である。
これが$L$に入ることは不可能。
残る$\nabla F(z)=0$の場合、偶側・奇側にそれぞれ二つ以上の零があり、
$z$は一つの偶座標と一つの奇座標が張る九本の座標直線の和に属する。
しかし$L$上のこれら二つの座標汎関数の制限は独立である。
もし従属なら$L$上で両方が同時に零となる点があり、
そこで$G=0$となるからである。
従って$L^\perp$はその九本の直線と交わらない。
以上で$X$は滑らかである。
一般の平面切断はBertiniの定理により滑らかな平面三次曲線で、
定理\ref{thm:ou-smooth-determinant-reconstruction}の仮定を満たす。
\end{proof}

\begin{theorem}[rank二の可換三次曲面としての認識]
\label{thm:ou-dense-rank-two-commutative}
定理\ref{thm:ou-dense-rank-two-smooth-surface}の仮定の下で、
$\Lambda$はそこで復元した滑らかな三次曲面$X$の
AOU\cite[(4.1), (4.2)]{AOU2024}の三ブロック九線束の
全自己準同型代数である（各ブロック内の順序を除く）。
従って標準noetherian OU helix比較を持つ曲面圏は
$\coh X$と同値である。これには内在的三次標識の存在を仮定しない。
\end{theorem}
\begin{proof}
$X\subset L^\perp\simeq\mathbb P^3$は滑らかである。
前証明から$L$には$G=0$を満たす境界点が存在しない。
従って$L^\perp$上の任意の二座標一次形式は独立であり、
偶側と奇側を含む任意の三座標一次形式も独立である：
従属関係は$L$内の台がその二つまたは三つの座標だけに
限られる元を与え、それは$G=0$の境界点となる。

各辺$(i,j)$を共有する偶奇置換$\pi,\pi'$に対して、
対応する二座標の共通零直線を$D_{ij}\subset\mathbb P^3$とする。
$F$の式によりこの直線は$X$上にある。
同じ置換に属する三辺の直線は、対応する座標平面の切断を成す。
上の一次形式の独立性から三直線は相異なる。
二つの異なる直線が同じ偶または奇置換を共有すれば、
同じ平面に入るので一点で交わり、滑らかな$X$上で交点数は一。
共有しなければ交点は$\nabla F=0$の座標直線に入ることになり、
前定理でこれを排除した。
従って元の$i,j$の添字では
\[
 D_{ij}^2=-1,\qquad
 D_{ij}\cdot D_{i'j'}=
 \begin{cases}1&i\ne i',\ j\ne j',\\0&i=i'\text{ または }j=j'\end{cases}
 \quad((i,j)\ne(i',j')).
\]
各置換の三直線の和は反標準因子$H_X=-K_X$。
同じ一辺を共有する二置換を比較して、
長方形の二辺ずつの和の線形同値を得る。
従って
\[
 U_0=\mathcal O_X,\quad V_j=\mathcal O_X(D_{0j}),\quad
 U_i=V_j(-D_{ij})
\]
は$j$に依らず定義される。
前節と同じ行積・列積のスカラー調整で、$A_{ij}$の置換積が
六座標切断そのものとなるように
$A_{ij}\in H^0(X,V_jU_i^{-1})$を選べる。

行列$\mathscr M_k=(c_{ijk}A_{ij})$は$X$上で行列式零である。
任意の閉点$p\in X$について$p$を通る滑らかな平面切断$E$が存在する。
実際、$p$を通る平面の線形系は$p$の外で基点を持たず、
Bertiniの定理と$p$での接平面を避ける条件による。
この$E$への制限は前節の滑らかな切断から復元した$U_i,V_j,A_{ij}$
と一致する（直線$D_{ij}$との交点が$p_{ij}$を与える）。
従って$\mathscr M_k|_E$は全点でrank二。
$p$は任意なので$\mathscr M_k$は$X$上で一定rank二である。
よって$W_k=(\ker\mathscr M_k)^\vee$は$X$上の線束であり、
核の双対から全ての非零射$B_{jk}:V_j\to W_k$を得る。
$W_k\cdot H_X=2$は$E$への制限から従い、
$W_k$の三つは$E$上で相異なるので$X$上でも相異なる。

次にこの九線束を古典的な六点爆発の表示で識別する。
$E_j=D_{0j}$ ($j=0,1,2$) は互いに素な$(-1)$曲線なので縮約できる。
縮約先$Y$は滑らかな次数六のdel Pezzo曲面である：
標準因子公式から$-K_Y$は全曲線に正の交点数を持ち、
自己交点六なのでampleである。
$i=1,2$について、$D_{ij}$の像は
$j$に依らない因子類$h_i$を持つ。
実際、その引戻しは
\[
 D_{ij}+\sum_{l\ne j}E_l=\sum_l E_l-U_i
\]
である（この式では線束を因子類で記す）。
上の交点数から$h_i^2=1$、$(-K_Y)h_i=3$、
$h_1h_2=2$を得る。
$h_i$は既約曲線で代表され自己交点が正なのでnefである。
次数六の滑らかなdel Pezzo曲面は平面の三非共線点の爆発である。
その表示で$h_i=aH-\sum_{\nu=4}^6 b_\nu E_\nu$と書くと
\[
 \sum b_\nu=3a-3,\qquad \sum b_\nu^2=a^2-1,\qquad b_\nu\geq0.
\]
Cauchy--Schwarzにより$1\leq a\leq2$である。
従ってこの類は$H$または$2H-E_4-E_5-E_6$であり、
どちらも平面への縮約の直線類である。
$h_1,h_2$は異なるので、平面表示を必要に応じCremona変換で替えて
$h_1=H$、$h_2=2H-\sum_{\nu=4}^6E_\nu$とできる。
$X$を六点の爆発と見たときの第一・第二ブロックは
（最初の三例外曲線の添字を$1,2,3$と改めて）
\[
 (0,-H+E_1+E_2+E_3,-2H+\sum_{\nu=1}^6 E_\nu),
 \qquad(E_1,E_2,E_3)
\]
となる。

最後に$W_k$の因子類を決める。
$B_{jk}$は非零で$H_X$に関する次数差が一なので、
$W_k-E_j$は有効な反標準次数一の因子、従って直線である。
$W_k=E_j+D$ ($D$は直線) と書ける。
$W_k^2$が負なら$E_j,D$が同じ直線または互いに素な直線であり、
$E_j\cdot W_k<0$となる。
この議論は$W_k^2<0$であれば各$j=1,2,3$に適用できる。
従って$W_k$の任意の有効代表は三本の$E_j$をすべて含み、
反標準次数二と矛盾する。ゆえに$W_k^2=0$であり、
$E_j\cdot W_k=0$ ($j=1,2,3$) となる。
（相異なる二直線の交点数は零または一である。）
二直線の和で自己交点零なら、その二直線との交点数も零で
その他の既約曲線との交点数は非負なので、$W_k$はnef。
従って$Y$上に降下するnef因子類であり、
その自己交点零、反標準次数二という条件は
\[
 \sum_{\nu=4}^6 b_\nu=3a-2,\qquad
 \sum_{\nu=4}^6 b_\nu^2=a^2,\qquad b_\nu\geq0
\]
を与える。Cauchy--Schwarzから
$6a^2-12a+4\leq0$、従って整数$a=1$。
三つの可能な類は$H-E_\nu$ ($\nu=4,5,6$) である。
三つの$W_k$は相異なるので、これら三類をすべて取る。

以上の九線束は順序を除いてAOUの(4.1)と一致する。
同論文のfull strong exceptional collectionとacyclic helixの結果を用いる。
隣接ブロックの非零射は$A_{ij},B_{jk}$であり、
$\mathscr M_k$の核の定義により元の関係を満たす。
その全自己準同型代数はAOU (4.2)の箙型であり、
隣接ブロックの18本の矢印で生成され、各$(i,k)$の関係空間は
一次元なので、元の$\Lambda$と同型である。
同じ$\Lambda$と逆Serreから導来helix全体が復元される。
滑らかなdel Pezzo曲面側では$s=-\otimes\mathcal O_X(-K_X)$で、
$-K_X$はampleなので、このhelixの商圏は$\coh X$に同値である
（周期をまとめた切断環によるArtin--Zhang再構成）。
標準noetherian OU helix比較の下で元の曲面圏は同じ商圏なので、
結論を得る。
\end{proof}

\begin{corollary}[全係数非零の可換locusの判定]
\label{cor:ou-dense-commutative-rank-criterion}
全係数非零でvery strong導来helixを持つ三ブロック関係テンソルに対し、
三つの巡回表示のいずれかで行列式線形式のrankが二であることと、
$\Lambda$が滑らかな可換三次曲面のAOU (4.1)の九線束型の
全自己準同型代数であることは同値である。
このとき全ての巡回表示のrankは二である。
それ以外の全係数非零型では、三つのrankは全て三で、
どの零平面も$T$に含まれない。
\end{corollary}
\begin{proof}
一方向は定理\ref{thm:ou-dense-rank-two-commutative}を
巡回表示に適用して得る。
ここで元のfoundationも同じ可換曲面上の九線束の
巡回操作を逆に戻したものとなる。
逆向きは、任意の巡回foundationの最初の二ブロックを
可換三次曲面$X$上で考える。
隣接ブロック間の九射はそれぞれ反標準次数一の非零切断なので
九本の直線を零因子とする。
同じ行または列の直線は相異なり
（異なればならないことはブロック内の対象の相異性による）、
その六置換積の線束は反標準束である。
最後の事実は最初のAOU foundationについて因子類から分かり、
巡回の場合もその九因子と$-K_X$を用いれば
隣接する二ブロックの和の差が$-K_X$であることが分かる。

$X\subset\mathbb P^3$の反標準埋込みにより、
六置換積は六つの非零一次形式の制限として表される。
その偶奇の三重積の差は$X$上で零である。
この差が$\mathbb P^3$上でも恒等的に零なら、
一意分解性で偶奇の三一次形式を比例する組に分けられる。
しかし偶置換と奇置換の零因子は、共有辺以外の二直線も
一致していなければならない。隣接二ブロックの九直線について
このことは起こらない。
実際、元のfoundationでは上記の因子類から九直線は相異なり、
巡回の場合も$V\to W$および$W\to sU$の九差は
AOU (4.1)の因子類から全て相異なる
（六点爆発の因子類の検算は下記コードを参照）。
従って三重積の差は$X$の三次方程式の非零定数倍である。
六一次形式が四次元空間を張らなければ、この三次方程式は
高々三つの一次形式だけに依存するconeとなって滑らかさに反する。
ゆえに六置換積は四次元の$H^0(X,-K_X)$を張る。
三つの$\ell_k$はこれら六置換積の線形関係なのでrankは高々二。
補題\ref{lem:ou-dense-determinant-rank}によりrank一は不可能、
従ってrank二となる。
すべての巡回表示について同じ議論を適用する。
\end{proof}

\begin{remark}
ここでの可換locusの判定は全係数非零かつvery strong導来helixを
持つ指定三ブロック箙型に限る。
三次syzygyは一般係数についても必要条件だが、それだけで
very strong性を保証しない。
また一般のOU三次箙型が全係数非零と仮定することはできない。
rank三の各零平面が特異な切断を与える型、
あるいは係数零を持つ型は引き続き調べる必要がある。
関連する整数計算は
\path{computations/verify_ou_determinant_cyclic.py}を参照。
\end{remark}
